\section{Limitations}
The released no-memory baseline uses the ordinal stream only. A matched shuffled no-memory run remains experiment debt. A per-task-restored rerun is the second prospective control. Until those runs are added, the 34.6--44.5\% and 14.9--17.5\% values remain task-exclusion sensitivity reductions. They do not constitute a causal mediation percentage for the full published gap.

The exposure detector uses a frozen rule set over explicit state channels. It omits mutations from scorer-failed writers when computing the 887-unit aggregate. It can also miss hidden service dependencies. The direct task-509 witness proves that the aggregate is a lower bound. An exposure label alone does not establish that a particular writer changed the exact answer required by a reader; the sentinel panel supplies that higher-precision evidence for three dynamic-reference tasks.

The empirical audit covers one WebArena release from one self-improvement study. The identification problem applies to stateful evaluation in general, while the reported prevalence values remain specific to this release. Multi-benchmark replication is therefore experiment debt for a broader empirical claim.

A substantial residual order gap remains after state-sensitive readers are excluded. Environment carryover therefore cannot account for the entire observed order effect. The paper's main empirical claim is material contamination of task-order measurement, with residual self-improvement sensitivity preserved as an active phenomenon.

The prospective contrast $\Delta_{\mathrm{mem}}-\Delta_{\mathrm{base}}$ assumes that the no-memory scaffold captures the environment-order contribution under the same execution stack. A memory method may interact with mutated state in ways the baseline does not reproduce. A factorial design with per-task restoration can resolve that interaction directly.

\section{Conclusion}
Task-order stress testing on a stateful benchmark is a joint intervention. The permutation changes the agent's experience history and the benchmark's service history. Released AWM trajectories contain a direct writer--reader ground-truth invalidation chain. The same mechanism reproduces perfectly across three dynamic-reference sentinels. Aggregate state-sensitive exclusion removes a material fraction of the reported order gap, especially for AWM.

These findings establish benchmark history as a first-class variable in longitudinal agent evaluation. The state-isolated order-stress protocol restores the intended scientific intervention by fixing service state while task order changes. Live-world variants can condition scoring on current service state. With benchmark history made explicit, task-order experiments can measure self-improvement history directly and support stronger claims about long-horizon agent evolution.
